%%%%%%%%%%%%%%%%%%%%%%%%%%%%%%%%%%%%%%%%%%%%%%%%%%%%%%%%%%%%%%%%%%%%%%
% Lớp 10 | Thống kê | Các số đặc trưng đo xu thế trung tâm của mẫu số liệu
% Dạng 0D6?3-3: Số trung vị 
% Số câu: 26  (sinh tự động bởi tach_cau.py; từ nay thầy sửa trực tiếp file này)
%%%%%%%%%%%%%%%%%%%%%%%%%%%%%%%%%%%%%%%%%%%%%%%%%%%%%%%%%%%%%%%%%%%%%%

%%%%%%%%%%%%%%%% TRẮC NGHIỆM 4 PHƯƠNG ÁN %%%%%%%%%%%%%%%%

% nguồn: MaHoaNH_TracNghiem_Lop10
\begin{ex}%[0D6N3-3]%[ma:f324dae8]
	Cho mẫu thống kê $28$; $16$; $13$; $18$; $12$; $28$; $13$; $19$ . Trung vị của mẫu  số liệu trên là bao nhiêu?
	\choice 
		 {$14$}
		{\True $17$} 
		 {$18$}
		 {$20$}
	\loigiai{ 
	 Mẫu thống kê trên có $ 8$ số liệu được sắp xếp theo thứ tự không giảm là $12$; $13$; $13$; $16$; $18$; $19$; $28$; $28$, nên trung vị của mẫu số liệu trên là $M_e=\dfrac{16+18}{2}=17$.}
\end{ex}

% nguồn: MaHoaNH_TracNghiem_Lop10
\begin{ex}%[0D6N3-3]%[ma:141c870c]
	Cho dãy số liệu thống kê  $32, 33, 36, 38, 39, 42, 46$. Khi đó trung vị là
	\choice 
		{$32$} 
		 {$46$}
		 {$37$}
		{\True $38$} 
	\loigiai{ 
		Mẫu số liệu sắp xếp theo thứ tự không giảm $32, 33, 36,  38, 39, 42, 46$.\\
		Cỡ mẫu $n=7$ nên trung vị của mẫu số liệu là $M_e=x_4=38$. 
	} 
\end{ex}

% nguồn: MaHoaNH_TracNghiem_Lop10
\begin{ex}%[0D6N3-3]%[ma:497c0bd7]
	Cho dãy số liệu thống kê $48$;  $36$; $33$; $38$; $32$; $48$; $42$; $33$; $39$. Khi đó số trung vị là
	 \choice
		 {$32$}
		 {$36$}
		{\True  $38$}
		 {$40$}
	\loigiai{ 
		Dãy số liệu thống kê được xếp thành dãy không giảm là $32$; $33$; $33$; $36$;  $38$; $39$; $42$; $48$; $48$.\\
		Ta có số trung  vị là $M_e=38$.
	 }
\end{ex}

% nguồn: MaHoaNH_TracNghiem_Lop10
\begin{ex}%[0D6N3-3]%[ma:29676ae2]
	Điểm thi học kì của một học sinh như sau: $4$;  $6$; $2$; $7$; $3$; $5$; $9$; $8$; $7$; $10$; $9$. Số trung bình và số trung vị lần lượt là
	 \choice
		 {$6{,}22$ và $7$}
		{$7$  và $6$}
		{\True $6{,}36$  và $7$}
		{$6$  và $6$}
	 \loigiai{
		Mẫu thống kê trên có $11$ số liệu được sắp xếp theo thứ tự không giảm là 
		\begin{center}
			$2$; $3$;  $4$; $5$; $6$; $7$; $7$; $8$; $9$; $9$; $10$.
		\end{center}
		Số trung bình là $\overline  {x}=\dfrac{2+3+4+5+6+2\cdot 7+8+2\cdot 9+10}{11}$ $=\dfrac{70}{11}\approx 6{,}36$.\\
		Trung vị của  mẫu số liệu trên là $M_e=7$.
	 }
\end{ex}

% nguồn: MaHoaNH_TracNghiem_Lop10
\begin{ex}%[0D6N3-3]%[ma:ff014114]
	Cân nặng (kg) của một nhóm học sinh nữ lớp $10$A cho bởi số liệu  sau:
	$$ 
	\begin{array}{llllllllll}
		38 & 39 & 44 & 44 & 46 & 47 & 47 & 48 & 52 & 55. 
	\end{array}
	$$ 
	Tìm  số trung vị của mẫu số liệu trên.
	\choice 
		{$44$} 
		{$46$} 
		{\True  $46{,}5$}
		 {$47$}
	\loigiai{ 
		Dãy trên đã sắp xếp không giảm, và có  $10$ số trong mẫu số liệu.\\
		 Số liệu thứ $5$ là $46$ và số liệu thứ $6$ là $47$.\\
	Vậy số trung vị  là $M_e=\dfrac{46+47}{2}=46{,}5$.}
\end{ex}

% nguồn: MaHoaNH_TracNghiem_Lop10
\begin{ex}%[0D6N3-3]%[ma:527d8ff6]
	Kết quả điểm kiểm tra môn Toán trong một kì thi của $200$ em học sinh  được trình bày ở bảng sau
	\begin{center}
		\begin{tabular}{|c|c|c|c|c|c|c|c|}
			\hline Điềm & $5$ & $6$ & $7$ & $8$ & $9$ & $10$  & Cộng\\
			 \hline Tần số & $10$ & $35$ & $38$ & $63$ & $42$ & $12$ & $200$\\
			\hline 
		\end{tabular}
	\end{center}
	Số trung vị của bản phân bố tần số nói trên  là
	 \choice
		 {\True $8$}
		 {$7$}
		 {$6$}
		 {$5$}
	\loigiai{ 
		 Số trung vị của bản phân bố tần số nói trên là $8$.
	} 
\end{ex}

% nguồn: MaHoaNH_TracNghiem_Lop10
\begin{ex}%[0D6N3-3]%[ma:1e8498be]
	Cho bảng số liệu thống kê chiều cao của một nhóm  học sinh như sau
	\begin{center}
		\begin{tabular}{|l|l|l|l|l|l|l|l|l|l|l|l|l|l|l|l|}
			\hline $150$ & $153$ & $153$ & $154$ & $154$ & $155$ & $160$ & $160$ & $162$ & $162$ & $163$ & $163$ & $163$ &  $165$ & $165$ & $167$\\
			\hline 
		\end{tabular}
	\end{center}
	Số trung vị của bảng  số liệu nói trên là
	 \choice
		{\True  $161$}
		{$153$} 
		 {$163$}
		{$156$} 
	\loigiai{ 
		Ta có trong bảng số liệu thống kê có tất cả  $16$ giá trị. Do đó số trung vị bằng trung bình cộng của hai số đứng thứ $8$ và $9$ trong bảng số liệu thống kê.\\
		Ta  có $M_e=\dfrac{160+162}{2}=161$.
	 }
\end{ex}

% nguồn: MaHoaNH_TracNghiem_Lop10
\begin{ex}%[0D6N3-3]%[ma:57d75064]
	Cho bảng số  liệu điểm bài kiểm tra môn toán của $20$ học sinh
	\begin{center}
		\begin{tabular}{|c|c|c|c|c|c|c|c|c|}
			\hline Điểm  & $4$ & $5$ & $6$ & $7$ & $8$ & $9$ & $10$ & Cộng\\
			\hline Số học sinh & $1$ & $2$ & $3$ & $4$ & $5$ & $4$ & $1$  & $20$\\
			\hline 
		\end{tabular}
	\end{center}
	Tìm số trung vị của bảng số  liệu trên.
	 \choice
		{$8$} 
		{$7$} 
		 {$7{,}3$}
		{\True  $7{,}5$}
	\loigiai{ 
		 Số trung vị của bảng số liệu có $20$ số là trung bình cộng của số thứ $10$ và số thứ $11$.\\
		Ta có số thứ $10$  là $7$; Số thứ $11$ là $8$.\\
	 Do đó $M_e=\dfrac{7+8}{2}=7{,}5$.}
\end{ex}

% nguồn: MaHoaNH_TracNghiem_Lop10
\begin{ex}%[0D6N3-3]%[ma:5271d16a]
	Cho bảng số liệu ghi lại điểm của $40$ học sinh trong bài kiểm tra $1$ tiết môn  Toán.
	\begin{center}
		\begin{tabular}{|c|c|c|c|c|c|c|c|c|c|}
			 \hline
			Điểm & $3$ & $4$ & $5$ & $6$  & $7$ & $8$ & $9$ & $10$ & Cộng\\
			\hline 
			Số học sinh & $2$ & $3$ & $7$ &  $18$ & $3$ & $2$ & $4$ & $1$ & $40$\\
			 \hline
		\end{tabular}
	\end{center}
	Số trung vị  là
	 \choice
		 {$5$}
		{\True $6$} 
		{$6{,}5$} 
		 {$7$}
	\loigiai{ 
		Số trung vị là  $M_e=\dfrac{x_{20}+x_{21}}{2}=\dfrac{6+6}{2}=6$.
	 }
\end{ex}

% nguồn: MaHoaNH_TracNghiem_Lop10
\begin{ex}%[0D6N3-3]%[ma:9228d2e9]
	Nhiệt độ $\left({}^{\circ}\mathrm{C}\right)$ trung bình các ngày trong tuần đầu tháng $3$ của  Nam Định là
	 $$
	\begin{array}{lllllll}
		20 & 19 & 22 &  19 & 24 & 28 & 25
	\end{array}
	 $$
	 Tìm số trung vị của mẫu số liệu trên.
	\choice 
		 {\True $22$}
		 {$23$}
		 {$24$}
		{$25$} 
	 \loigiai{
		Sắp xếp theo  thứ tự không giảm:
		$\begin{array}{ccccccc}19 & 19 & 20 & 22 & 24 & 25 & 28.
	\end{array}$
	 Mẫu số liệu trên có $7$ số.\\
	Số trung vị là số đứng ở vị trí thứ $\dfrac{7+1}{2}=4$.\\ 
Vậy  số trung vị là $M_e=22$}
\end{ex}

% nguồn: MaHoaNH_TracNghiem_Lop10
\begin{ex}%[0D6N3-3]%[ma:2c81adb3]
	Số điểm  kiểm tra $11$ môn của một nhóm gồm $11$ học sinh được cho trong bảng sau:
	\begin{center}
		\begin{tabular}{|c|c|c|c|c|c|c|c|}
			\hline 
			Điểm & $4$ & $5$ & $7$ & $8$  & $9$ & $10$ &\\
			 \hline
			Tần số & $2$ & $1$ & $2$ &  $3$ & $1$ & $2$ & $N=11$\\
			 \hline
		\end{tabular}
	\end{center}
	Số trung vị của mẫu số liệu trên  là
	\choice 
		 {$7$}
		 {\True $7{,}5$}
		 {$8$}
		 {$8{,}5$}
	\loigiai{ 
		Số trung vị là $M_e=\dfrac{7+8}{2}=7{,}5$. 
	} 
\end{ex}

% nguồn: MaHoaNH_TracNghiem_Lop10
\begin{ex}%[0D6H3-3]%[ma:20d4729f]
	Với $x$ nguyên dương, cho mẫu số liệu sau (đã sắp xếp theo thứ  tự)
	$$ 
	\begin{array}{lllllllll}
		 4 & 9 & 10 & 4x-3 & 4x-1 & 4x & 22 & 23 & 25
	\end{array}
	$$ 
	Tìm $x$, biết rằng số trung vị trong mẫu  số liệu trên bằng $15$.
	\choice 
		 {\True $x=4$}
		{$x=12$} 
		{$x=15$} 
		{$x=14$} 
	\loigiai{ 
		Ta có $M_e=4x-1=15 \Rightarrow x=4$. 
	} 
\end{ex}

% nguồn: MaHoaNH_TracNghiem_Lop10
\begin{ex}%[0D6H3-3]%[ma:5dbd16f9]
	Số liệu sau đây cho ta lãi hàng tháng của một cửa hàng trong năm $2023$. Đơn vị là triệu  đồng
	\begin{center}
		\begin{tabular}{|c|c|c|c|c|c|c|c|c|c|c|c|c|}
			\hline Tháng & 1 & 2 & 3 & 4 & 5 & 6 & 7 & 8 & 9 & 10 & 11 & 12 \\ 
			\hline Lãi & 10 & 12 & 18 & 15 & 22 & 20 & 15 &  13 & 18 & 17 & 20 & 22 \\
			\hline 
		\end{tabular}
	\end{center}
	Tìm số trung vị của mẫu  số liệu trên.
	 \choice
		{$15$} 
		 {\True $17{,}5$}
		{$20$} 
		{$21$} 
	 \loigiai{
		Sắp xếp tiền lãi theo thứ tự không giảm 
		\begin{center}
			\begin{tabular}{|c|c|c|c|c|c|c|c|c|c|c|c|}
				\hline 10 & 12 & 13 & 15 & 15 & 17 & 18 & 18 & 20 & 20 & 22 & 22 \\ 
				\hline 
			\end{tabular}
		\end{center}
		Mẫu số liệu trên có $12$ số. Số thứ $6$ và số thứ $7$ lần  lượt là $17$ và $18$.\\
	 Vậy số trung vị là $M_e=\dfrac{17+18}{2}=17{,}5$.}
\end{ex}

% nguồn: MaHoaNH_TracNghiem_Lop10
\begin{ex}%[0D6H3-3]%[ma:0d057b4e]
	Khi kiểm tra ngẫu nhiên một số công nhân trong một xí  nghiệp, người ta thống kê lại độ tuổi của họ ở bảng sau
	\begin{center}
		\begin{tikzpicture}[xscale=1.7,yscale=.7]
				\begin{scope}[shift={(-.5,.5)}]
					\fill[cyan!10]  (-2.2,0) rectangle (0,-2);
					\draw (6,-2) grid  (0,0) rectangle (-2.2,-2)
						 (0,-1)--(-2.2,-1);
				\end{scope}
				 \path
					(-2.5,0)  node[right]{Tuổi}
					(-2.5,-1) node[right]{Số công nhân} 
					(0,0)  node{25}
					(0,-1) node{4} 
					 (1,0) node{26}
					 (1,-1) node{9}
					(2,0)  node{27}
					(2,-1)  node{8}
					 (3,0) node{29}
					 (3,-1) node{3}
					 (4,0) node{31}
					 (4,-1) node{1}
					(5,0)  node{34}
					(5,-1) node{1}; 
			\end{tikzpicture}
	\end{center}
	Tìm trung vị của mẫu số  liệu trên.
	\choice 
		{$27$} 
		 {$26$}
		{$25{,}6$} 
		{\True  $26{,}5$}
	 \loigiai{
	Kí hiệu $x_1$, $x_2$, $\ldots$,  $x_n$ là các giá trị của mẫu số liệu trên, với $x_k\le x_{k+1}$.\\
	Ta có $n=4+9+8+3+1+1= 26$. Vậy số trung vị cần tìm  là $Q_2= \dfrac{x_{13}+x_{14}}{2}= \dfrac{26+27}{2}= 26{,}5$.
	 }
\end{ex}

% nguồn: Ok_TracNghiemLop11_Lan1
\begin{ex}%[0D6B3-3]%[ma:9fe83b02]
	 Gọi $M=\cos \left(a+b\right)\cdot\cos \left(a-b\right)+\sin \left(a+b\right)\cdot\sin \left(a-b\right)$ thì
	\choice 
		 {$M=\sin4b$}
		 {\True $M=1-2\sin^2b$}
		 {$M=1+2\sin^2b$}
		{$M=\cos 4b$} 
	 \loigiai{
		Ta  có
		\begin{eqnarray*}
			M&=&\cos \left(a+b\right)\cdot\cos \left(a-b\right)+\sin \left(a+b\right)\cdot\sin \left(a-b\right)\\ 
			&=&\cos  \left(\left(a+b\right)-(a-b)\right)=\cos 2b=1-2\sin^2b
		\end{eqnarray*}
	 }
\end{ex}

% nguồn: Ok_TracNghiemLop11_Lan1
\begin{ex}%[0D6B3-3]%[ma:5a9a11e7]
	Biểu thức $2\sin  \left(\dfrac{\pi}{4}+\alpha \right)\sin \left(\dfrac{\pi}{4}-\alpha \right)$ đồng nhất với biểu thức nào dưới đây ?
	 \choice
		 {$\sin 2\alpha$}
		{\True $\cos  2\alpha$}
		 {$\sin \alpha$}
		{$\cos  \alpha$}
	\loigiai{ 
	 Ta có $2\sin \left(\dfrac{\pi}{4}+\alpha \right)\sin \left(\dfrac{\pi}{4}-\alpha \right)=2 \cdot \dfrac{1}{2}\left[\cos 2\alpha -\cos \dfrac{\pi}{2}\right]=\cos 2\alpha$.}
\end{ex}


%%%%%%%%%%%%%%%% ĐÚNG/SAI %%%%%%%%%%%%%%%%

% nguồn: MaHoaNH_DungSai_Lop10
\begin{ex}%[0D6N3-3]%[ma:5c912b09]
	Bảng số liệu sau cho biết mức lương hàng năm của các cán bộ và nhân  viên trong một công ty (đơn vị: nghìn đồng)
	\begin{center}
		\begin{tabular}{|c|c|c|c|}
			 \hline
			$20910$  & $76000$ & $20350$ & $20060$\\
			 \hline
			$21410$ & $20110$ &  $21410$ & $21360$\\
			 \hline
			$20350$ & $21130$ & $20960$ &  $125000$\\
			 \hline
		\end{tabular}
	\end{center}
	Xét tính đúng sai của  các mệnh đề sau
	 \choiceTF
		{\True Mức  lương trung bình các cán bộ nhân viên là $\overline {x}=34087{,}5$ (nghìn đồng)}
		{Mức lương  lớn nhất là $76000$}
		{\True Số trung vị  là $21045$ (nghìn đồng)}
		{Có thể lấy mức lương bình quân làm giá trị đại diện} 
	\loigiai{ 
		 Sắp theo thứ tự không giảm bảng lương ta được\\
		$20060$; $20110$; $20350$;  $ 20350$; $20910$; $20960$; $21130$; $21360$; $21410$; $21410$; $76000$; $125000$, \quad $(n=12)$.\\
		Mức lương trung bình các cán bộ  nhân viên là\\
		$\overline {x}=\dfrac{20910+76000+\ldots+125000}{12}=\dfrac{68175}{2}=34087{,}5$ (nghìn  đồng).\\
		Số trung vị là $\dfrac{20960+21130}{2}=21045$ (nghìn đồng).\\ 
		Ý nghĩa: Số trung vị đại diện cho mức lương trung bình của cán bộ nhân  viên của công ty, vì trong trường hợp này chênh lệch giữa các số liệu quá lớn nên không thể lấy mức lương bình quân làm giá trị đại diện.\\
		 Do đó
		\begin{itemchoice}
			\itemch Đúng. 
			\itemch  Sai.
			\itemch  Đúng.
			 \itemch Sai.
		\end{itemchoice}
	} 
\end{ex}

% nguồn: MaHoaNH_DungSai_Lop10
\begin{ex}%[0D6N3-3]%[ma:906aba09]
	Tập đoàn $X$ có $24$ công ty. Thống kê cuối năm cho biết doanh thu (đơn vị triệu đồng) của $24$ công  ty con như sau
	\begin{center}
		\begin{tabular}{|c|c|c|c|c|c|c|c|}
			 \hline
			$35\,432$ & $14\,215$ & $24\,436$ & $13\,978$ &  $45\,713$ & $16\,323$ & $37\,488$ & $13\,458$\\
			\hline 
			$57\,754$ & $53\,345$ & $80\,234$ & $117\,245$ &  $74\,506$ & $86\,851$ & $47\,678$ & $611\,298$\\
			 \hline
			$19\,397$ & $48\,644$ & $8\,324$ & $9\,599$  & $94\,338$ & $45\,390$ & $37\,492$ & $811\,854$\\
			\hline 
		\end{tabular}
	\end{center}
	Xét tính đúng sai  của các mệnh đề sau
	 \choiceTF
		{Doanh thu thấp nhất là  $9\,599$}
		{\True Doanh  thu lớn nhất là $811\,854$}
		{\True Số  trung bình của mẫu số liệu trên khoảng $100\,208$}
		{\True Số trung vị là $45\,551{,}5$} 
	\loigiai{ 
		Sắp xếp  lại mẫu số liệu trên theo thứ tự tăng dần từ trái qua phải, từ trên xuống
		\begin{center}
			\begin{tabular}{|c|c|c|c|c|c|c|c|}
				 \hline
				$8\,324$ & $9\,599$ & $13\,458$ & $13\,978$ & $14\,215$ & $16\,323$ & $19\,397$ &  $24\,436$\\
				\hline 
				$35\,432$ & $37\,488$ & $37\,492$ & $45\,390$ & $45\,713$ & $47\,678$ &  $48\,644$ & $53\,345$\\
				 \hline
				$57\,754$ & $74\,506$ & $80\,234$ & $86\,851$ & $94\,338$  & $117\,245$ & $611\,298$ & $811\,854$\\
				 \hline
			\end{tabular}
		\end{center}
		Số trung bình của mẫu  số liệu trên khoảng $100\,208$.\\
		Số trung vị là $45\,551{,}5$.\\ 
		Vì giá trị trung bình lớn hơn giá trị trung vị  rất nhiều, điều đó thể hiện mẫu có
		một số giá trị bất thường. Vì vậy, số trung vị làm đại  diện mẫu sẽ tốt hơn.
		Do đó 
		\begin{itemchoice}
			\itemch  Sai.
			\itemch Đúng. 
			\itemch  Đúng.
			\itemch  Đúng.
		\end{itemchoice}
	} 
\end{ex}

% nguồn: MaHoaNH_DungSai_Lop10
\begin{ex}%[0D6N3-3]%[ma:a13b994a]
	Thống kê chiều cao (đơn vị cm) của nhóm $15$ bạn nam lớp $10$ cho kết quả  như sau\\
	\begin{center}
		\begin{tabular}{|c|c|c|c|c|c|c|c|c|c|c|c|c|c|c|}
			\hline 
			$162$ &  $157$ & $170$ & $165$ & $166$ & $157$ & $159$ & $164$ & $172$ & $155$ & $156$ & $156$ & $180$ & $165$ & $155$\\
			 \hline
		\end{tabular}
	\end{center}
	Xét tính  đúng sai của mệnh đề
	 \choiceTF
		{Chiều cao thấp  nhất là $156$}
		 {\True $Q_2=162$}
		{$Q_1=157$} 
		 {$Q_3=170$}
	 \loigiai{
		 Sắp xếp số liệu theo thứ tự tăng dần ta được
		\begin{center}
			\begin{tabular}{|c|c|c|c|c|c|c|c|c|c|c|c|c|c|c|}
				 \hline
				$155$ & $155$ & $156$ & $156$ & $157$ & $157$ & $159$ & $162$ & $164$ & $165$ & $165$ &  $166$ & $170$ & $172$ & $180$\\
				\hline 
			\end{tabular}
		\end{center}
		Vì có $15$ giá trị nên số trung  vị là số ở vị trí thứ $8$: ${Q_2}=162$.\\
		Nửa số liệu bên trái  $Q_2$ là
		\begin{center}
			\begin{tabular}{|c|c|c|c|c|c|c|}
				 \hline
				$155$ & $155$ & $156$ & $156$ & $157$ & $157$ &  $159$\\
				\hline 
			\end{tabular}
		\end{center}
		 Ta tìm được trung vị $Q_1=156$.\\
		Nửa số liệu bên  phải $Q_2$ là
		\begin{center}
			\begin{tabular}{|c|c|c|c|c|c|c|}
				 \hline
				$164$ & $165$ & $165$ &  $166$ & $170$ & $172$ & $180$\\
				 \hline
			\end{tabular}
		\end{center}
		Ta tìm được trung  vị $Q_3=166$.
		Do  đó
		\begin{itemchoice}
			\itemch  Sai.
			\itemch  Đúng.
			 \itemch Sai.
			\itemch Sai. 
		\end{itemchoice}
	} 
\end{ex}


%%%%%%%%%%%%%%%% TRẢ LỜI NGẮN %%%%%%%%%%%%%%%%

% nguồn: MaHoaNH_TraLoiNgan_Lop10
\begin{bt}%[0D6N3-3]%[ma:029a9d88]
	Chiều dài (đơn  vị feet) của $7$ con cá voi trưởng thành được cho như sau $48$; $53$; $51$; $31$; $53$; $112$; $52$. Tìm trung vị của mẫu số liệu trên.
	\shortans{52} 
	\loigiai{ 
		Sắp xếp lại mẫu số liệu không giảm như sau $31$; $48$; $51$; $52$; $53$;  $53$; $112$.\\
		Do đó trung vị của mẫu là $M_e=52$. 
	} 
\end{bt}

% nguồn: MaHoaNH_TraLoiNgan_Lop10
\begin{bt}%[0D6H3-3]%[ma:97999006]
	Cho  mẫu số liệu: $2\quad 3\quad 7\quad 4\quad 5\quad 3\quad 4\quad 4$. Tính trung vị của mẫu số liệu trên.
	 \shortans{$4$}
	 \loigiai{
		Sắp xếp mẫu: $2\quad  3\quad 3\quad 4\quad 4\quad 4\quad 5\quad 7$.\\
		Kích thước mẫu là $8$.\\ 
		Khi đó trung vị của mẫu là trung bình cộng  của giá trị thứ $4$ và $5$ là $M_e=\dfrac{1}{2}\cdot(4+4)=4$.
	 }
\end{bt}

% nguồn: MaHoaNH_TraLoiNgan_Lop10
\begin{bt}%[0D6H3-3]%[ma:37c634ee]
	Bảng sau cho biết thời gian chạy cự li $100$\,m của các bạn trong lớp (đơn  vị giây)
	\begin{center}
		\begin{tabular}{|l|c|c|c|c|c|}
			\hline Thời gian & $12$ &  $13$ & $14$ & $15$ & $16$ \\
			\hline Số  bạn & $4$ & $7$ & $3$ & $18$ & $8$ \\
			\hline 
		\end{tabular}
	\end{center}
	Số tứ phân vị $Q_3$ của mẫu số  liệu trên bằng bao nhiêu?
	\shortans{15} 
	 \loigiai{
		Số bạn học sinh  trong lớp là $n=4+7+3+18+8=40$ (bạn).\\
	Tứ vị phân  thứ ba là $Q_3=\dfrac{15+15}{2}=15$.}
\end{bt}

% nguồn: MaHoaNH_TraLoiNgan_Lop10
\begin{bt}%[0D6H3-3]%[ma:40346f27]
	Một cửa hàng giày thể  thao đã thống kê cỡ giày của $20$ khách hàng nữ được chọn ngẫu nhiên cho kết quả như sau:
	 $$
	35\quad 37\quad 39\quad 41\quad 38\quad 40\quad 40\quad  37\quad 39\quad 38\quad 38\quad 36\quad 37\quad 42\quad 38\quad 35\quad 38\quad 36\quad 38\quad 35.
	$$ 
	Số trung vị của mẫu số  trên bằng bao nhiêu?
	\shortans{38} 
	\loigiai{ 
		Sắp xếp các giá trị theo  thứ tự không giảm:
		 $$
		35\quad 35\quad 35\quad 36\quad 36\quad 37\quad  37\quad 37\quad 38\quad 38\quad 38\quad 38\quad 38\quad 38\quad 39\quad 39\quad 40\quad 40\quad 41\quad 42 .
		$$ 
		Vì $n=20$ là số chẵn nên trung vị là trung bình cộng của hai giá  trị chính giữa
		$$ 
	M_e=\dfrac{38+38}{2}=38.$$} 
\end{bt}

% nguồn: Ok_TraLoiNgan_Lop11_Lan1
\begin{ex}[c]%[0D6B3-3]%[ma:0b59873d]
	Tính $M=\cos a+\cos \left(a+120^0\right)+\cos \left(a-120^0\right)$. 
	\shortans{0} 
	 \loigiai{
	Áp dụng công thức biến đổi tổng thành tích được \[M=\cos a+2\cos a.\cos 120^\circ =\cos a+2\cos a.\left(-\dfrac{1}{2}\right)=\cos  a-\cos a=0\]}
\end{ex}

% nguồn: Ok_TraLoiNgan_Lop11_Lan1
\begin{ex}%[0D6B3-3]%[ma:ef391229]
	Cho hai góc nhọn $a$ và $b$. Biết $\cos a=\dfrac{1}{3}$, $\cos b=\dfrac{1}{4}$. Giá trị $\cos\left(a+b\right)\cdot\cos\left(a-b\right)$ có dạng $-\dfrac{a}{b}$ với $(a,  b) =1$.Khi đó $a + b$
	\shortans{263} 
	\loigiai{ 
	Ta có $$\cos\left(a+b\right)\cdot\cos\left(a-b\right)=\dfrac{1}{2}\left(\cos 2a+\cos 2b\right)=\cos  ^2a+\cos ^2b-1=\left(\dfrac{1}{3}\right)^2+\left(\dfrac{1}{4}\right)^2-1=-\dfrac{119}{144}.$$}
\end{ex}


%%%%%%%%%%%%%%%% TỰ LUẬN %%%%%%%%%%%%%%%%

% nguồn: MaHoaNH_TuLuan_Lop10
\begin{bt}%[0D6H3-3]%[0D6H3-2]%[ma:7cd4b663]
	Cho bảng số liệu ghi lại điểm của $40$ học  sinh trong bài kiểm tra giữa kỳ I môn Toán như sau
	\begin{center}
		\begin{tabular}{|c|c|c|c|c|c|c|c|c|c|}
			\hline Điểm & $3$ & $4$ & $5$ & $6$ & $7$ & $8$ & $9$ & $10$ & Cộng\\ 
			\hline  Số học sinh & $2$ & $3$ & $7$ & $18$ & $3$ & $2$ & $4$ & $1$ & $40$\\
			\hline 
		\end{tabular}
	\end{center}
	Tìm số trung bình và trung vị của mẫu số liệu trên. 
	 \loigiai{
		Số trung bình  của mẫu số liệu trên
		\allowdisplaybreaks 
		\begin{eqnarray*}
			\overline{x}=\dfrac{3\cdot 2+4\cdot  3+5\cdot 7+6\cdot 18+7\cdot 3+8\cdot 2+9\cdot 4+10\cdot 1}{40}=6{,}1.
		\end{eqnarray*}
		Xếp các giá trị theo chiều  không giảm. Theo công thức trung vị đối với cỡ mẫu $n=40$ chẵn thì ta có số đứng vị trí $20$ là $6$ và số đứng vị trí $21$ là $6$.\\
		Vậy số trung vị $M_e=6$ 
	 }
\end{bt}
